\section{格点微扰论}
\label{sec:per}

为了对下一节给出的非微扰结果有所洞察，我们在连续统与格点上、于朗道规范中完成了相应的一圈微扰计算。
该计算是文献 \cite{newi} 中计算的推广。

标准的一圈修正计算，对应于在圈动量上取固定紫外截断 $\vert p_i \vert \le \pi/a$ 来计算重正化常数，
并在最终结果中略去 $O(a)$ 项。在此极限下得到的重正化常数对夸克质量 $m_q$ 与外动量的依赖仅通过
对数发散项体现，这些项与算符的反常量纲相关。由于在一圈处这些项是普适的——即在格点与任何连续统正规化中
都相同——联系格点重正化算符与连续统重正化算符所需的重正化常数不依赖于夸克质量和外动量（但依赖于格点间距
与连续统重正化标度 $\mu$）。我们称这一极限为``标准微扰论''（STP）。然而非微扰计算是在有限大小的格点上、
在固定的格点间距与夸克质量值下进行的：因此无法使 $O(a)$ 项任意小。特别地，由于我们想研究重正化常数的
$\mu$ 依赖性，重要的是知道在哪些 $\mu$（及 $m_q$）值处，在有限格点、固定格点间距下得到的微扰结果会偏离 STP
的结果。为紧贴非微扰计算，微扰结果也在 $16^3 \times 32$ 的格点上求出。这样，最终结果中保留了所有
$O(\alpha_s a)$ 或 $O(a^2)$ 或更高阶的项（我们称此流程为``离散微扰论''，DPT）。圈积分在 SPT 与 DPT 中都是
数值计算的。

非微扰结果是通过对在离壳夸克态之间、朗道规范中计算的适当截腿格林函数求迹得到的，参见式 (\ref{eq:proj})。
在微扰论中，这一流程与标准流程不同：标准流程是从正比于所考虑算符树图级矩阵元的项中辨认出修正
\cite{buras,ciu}。差别源于一些有限项：在标准流程中它们被当作算符矩阵元的一部分，而在投影算符方案中则被当作
一圈修正的一部分。在一圈阶，这些项在连续统与格点上相同（在 $a
\to 0$ 极限），并且在计算连续统与格点重正化常数之差时相消。
\begin{figure}[htbp]
\centering
\includegraphics[width=0.62\textwidth]{images/feynman2.pdf}
\caption{Clover 作用量的顶点图。第二、三幅图对应于算符转动部分的贡献。}
\label{fig:vertex}
\end{figure}

说得更具体些，考虑图 \ref{fig:vertex} 中的图，在朴素维数正规化（NDR）、费曼规范下计算：
\beq \Lambda_{\gamma_\mu}(\frac{p}{\mu})=
<p|\bar \psi \gamma_\mu \psi|p> =
 \Bigl[ 1+\frac{\alpha_s}{4 \pi} C_F \Bigl(\frac{p^2}{\mu^2}\Bigr)^{-\epsilon}
\Bigl( (\frac{1}{\hat \epsilon}+1) \gamma_\mu-2\frac{\pslash p_\mu}{p^2}
\Bigr) \Bigr]   ,\label{eq:gcont}\eeq
其中 $C_F=(N^2-1)/2N$，$1/\hat \epsilon=1/\epsilon
-\ln 4 \pi - \gamma_E$，$p$ 是外动量。通常把正比于 $(1/\hat \epsilon +1)$ 的项，即
$\gamma_\mu$ 的系数，作为裸算符的一圈贡献。
而使用投影算符（见式 (\ref{eq:proj})）时，式 (\ref{eq:gcont}) 的最后一项也有贡献，于是因子变为
$(1/\hat \epsilon+1/2)$。表 \ref{tab:cont} 给出了用投影算符、在 't Hooft--Veltman（HV）维数重正化方案
\cite{hv} 下、对本文所考虑的两夸克算符得到的一圈修正。NDR 正规化与 HV 方案的差别在于 $\gamma_5$
反交换性质的定义；细节例如见 \cite{ciu}。在 NDR 中，轴矢流与赝标密度的修正分别与矢量流和标量密度的
修正一致。把胶子传播子写成
$G(k)=\left(-\delta_{\mu\nu}k^2+(1-\lambda) k_\mu k_\nu \right)/k^4$，
我们把一圈修正中正比于 $(1-\lambda)$ 的贡献称为``朗道''部分。
\begin{table}[htbp]
\begin{center}
\begin{tabular}{||c|c|c||}
%
\hline
%
\textrm{算符}
&费曼&朗道   \\ \hline \hline
%
%
$\gamma_\mu$ &$1/\hat \epsilon+ 1/2$ &$-1/\hat
\epsilon-1/2$  \\
\hline
$\gamma_\mu\gamma_5$&$1/\hat \epsilon+5/2$ &$-1/\hat
\epsilon-1/2$     \\ \hline
$\hat 1$  &$4/\hat \epsilon+6$ &$-1/\hat
\epsilon-2$ \\ \hline
%
$\gamma_5$  &$4/\hat \epsilon+10$ &$-1/\hat
\epsilon-2$ \\ \hline
%
$\Sigma$  &$-1/\hat \epsilon-1$ &$1/\hat
\epsilon+1$ \\  \hline
\end{tabular}
\caption[]{ \itshape 图 \ref{fig:vertex} 的一圈图对不同两夸克算符在 HV 方案下的计算结果。同时给出图
\ref{fig:self} 自能图的计算结果。夸克质量取为零，逆夸克传播子写为 $i \pslash \Sigma(p^2)$，并省略了因子
$(\alpha_s/4 \pi) C_F \Bigl(p^2/\mu^2\Bigr)^{-\epsilon}$。}
\label{tab:cont}
\end{center}
\end{table}
\begin{figure}[htbp]
\centering
\includegraphics[width=0.62\textwidth]{images/feynman3.pdf}
\caption{自能图。}
\label{fig:self}
\end{figure}

采用 $\overline{{\rm MS}}$ 方案，减除极点后，图 \ref{fig:vertex} 的不可约顶点形如\footnote{$
\Gamma_O(p/\mu)$ 由 $\Lambda_O(p/\mu)$ 经式 (\ref{eq:proj}) 所述的投影算符得到。}
 \beq \Gamma_O(p/\mu)=
  1+\frac{\alpha_s}{4 \pi} C_F
\Bigl(-\gamma^{(O)} \ln (p^2/\mu^2)+C_O^{\overline{{\rm MS}}}
 \Bigr). \label{eq:gms}\eeq
{}从表 \ref{tab:cont} 可得：费曼规范下 $C_{\gamma_\mu}^{\overline{{\rm
MS}}} =1/2$，$C_{\hat 1}^{\overline{{\rm MS}}}
 =6$ 等；而对朗道（纵向）项有
$C_{\gamma_\mu}^{\overline{{\rm MS}}} =-1/2$，$C_{\hat 1}^{\overline{{\rm MS}}}
 =-2$ 等。

在格点上可得类似表达式：
\beq \Gamma_O(pa)=
  1+\frac{\alpha_s}{4 \pi} C_F
\Bigl(-\gamma^{(O)} \ln (p^2 a^2)+C_O^{{\rm Latt}}
 \Bigr). \label{eq:gl}\eeq
$C_O^{{\rm Latt}}$ 对费曼与朗道修正的数值见表 \ref{tab:latt}。
\begin{table}[htbp]
\begin{center}
\begin{tabular}{||c|c|c|c||}
%
\hline
%
\textrm{算符}
&费曼&转动后的费曼&朗道   \\ \hline \hline
%
%
$\gamma_\mu$ &$6.62$ &$-3.52$ &$-4.29$ \\
\hline
$\gamma_\mu\gamma_5$&$5.09$ &
$-11.68$ &$-4.29$ \\ \hline
%
$\hat 1$ &$16.11$&$-12.00$&$-5.79$ \\ \hline
$\gamma_5$  &$19.18$ &$4.31$ &$-5.79$\\ \hline
%
$\Sigma$  &$8.21$ &$--$&$4.79$ \\  \hline
\end{tabular}
\caption[]{ \itshape Clover 作用量下不同两夸克算符及自能的 $C_O^{{\rm Latt}}$
（$C_\Sigma$）数值。夸克质量取为零。}
\label{tab:latt}
\end{center}
\end{table}

表 \ref{tab:latt} 中我们把两费米子算符转动部分（见式 (\ref{eq:oi}) 与 (\ref{eq:rotaz})）的贡献记为
``转动后的费曼''。算符转动部分的朗道贡献为零。这是由如下事实要求的：联系格点算符与连续统算符的重正化
常数应是规范不变的。

遵循文献 \cite{newi}--\cite{2ferm}，我们写出
\beq O(\mu)=\left(1+\frac{\as}{4 \pi} C_F \left(-\gamma^{(0)}\ln (\mu a)^2
+\Delta_O \right) \right) O(a) ,\eeq
其中 $\Delta_O$ 可由表 \ref{tab:cont} 与 \ref{tab:latt} 计算：
\beq \Delta_O=C_O^{\overline{{\rm MS}}}-C_O^{{\rm Latt}} +
 C_\Sigma^{\overline{{\rm MS}}}-C_\Sigma^{{\rm Latt}} .\eeq
在连续统与格点一圈修正之差中，如预期的那样，所有朗道部分都相消。
$\Delta_{\gamma_\mu},
\Delta_{\gamma_\mu \gamma_5},\cdots$ 的结果均与文献 \cite{newi} 一致。

在格点上，对给定的夸克质量与外动量，存在 $a$ 阶的项，它们在大 $m_q$ 与大 $p$ 处可能变得重要。由于我们的
方法必须在较大的动量值处施加重正化条件，至少在微扰论中监测 $O(a)$ 效应变得重要的质量与动量取值是重要的。
DPT 中格林函数的形式类似于 (\ref{eq:gl})：
\beq \Gamma_O(pa)=
  1+\frac{\alpha_s}{4 \pi} C_F
\Bigl(-\gamma^{(O)} \ln (p^2 a^2)+C_O^{{\rm Latt}}(pa,m_q a)
 \Bigr), \label{eq:gld}\eeq
只是``常数'' $C_O^{{\rm Latt}}$ 现在依赖 $O(a)$ 项。因此在 DPT 的一圈修正计算中包含了所有依赖截断的项。
在有限体积上，无质量粒子零动量的传播子需要调节。费曼图的结果依赖调节方式，但这种依赖随体积增大而消失。
在下文给出的结果中，我们把夸克与胶子传播子中的零动量项设为零。我们观察到低 $\mu^2$ 处 SPT 与 DPT 结果之间
的小差别即源于这一选择。在我们进行模拟所用的夸克质量值处，质量修正是可忽略的。

出现在式 (\ref{eq:gl}) 与 (\ref{eq:gld}) 中的耦合常数是裸格点耦合。文献 \cite{lepage} 中论证了可以通过用
以物理量定义的 ``boosted'' 耦合 $\alpha_s^V$ 重组展开来优化微扰级数的收敛性。在与下面非微扰结果的比较中，
我们遵循文献 \cite{lepage} 的建议，使用有效耦合计算微扰修正：
\beq \alpha_s^V= \frac{1}{\langle \frac{1}{3}{\mathrm{Tr}} U_{plaq}\rangle}
 \alpha_s^{LATT} \simeq  1.68\,\, \alpha_s^{LATT} ,\eeq
其中 $\alpha_s^{LATT}=g_0^2/4 \pi$，$\langle \frac{1}{3}{\mathrm{Tr}} U_{plaq}\rangle$ 是方格（plaquette）的
期望值。我们称此展开为``Boosted 离散微扰论''（BDPT）；类似地，``Boosted 标准微扰论''记作 BSPT。

BDPT 中重正化常数的计算结果将在下一节与非微扰结果比较。
